\section*{4 From Hypergraphs to Parity Samplers}
Fix some $n,t,d\in\mathbb N$ and a $d$-regular3-uniform hypergraph $H=(V,E_H)$ over $n$ vertices. We will construct a parity sampler $\mathbf W_{H,t}\sim[n]^t$ from $H$ and show that when $H$ is a good spectral expander, $\mathbf W$ is a good parity sampler.
\paragraph{The construction.}
For each $v\in V$, there are $d$ edges $(v_1,v_2,v_3)\in E_H$ with $v_1=v$. For any $i\in[d]$, let $e_H(v,i)$ be the $i$th such edge (i.e., $e_H(v,i)_1=v$). Without loss of generality, we assume the vertices are each labeled with a unique integer between1 and$n$ (i.e $V=[n]$).
To sample $\mathbf w\sim\mathbf W_{H,t}$, independently sample a starting vertex $\mathbf v_0\in[n]$ and edge labels $\mathbf i_1,\ldots,\mathbf i_t\sim[d]$ uniformly. $\mathbf w$ will be a deterministic function of $\mathbf v_0$ and $\mathbf i_1,\ldots,\mathbf i_t$ computed as follows. For each $j=1,\ldots,t$,
\begin{enumerate}
\item Let $\mathbf e_j=e_H(\mathbf v_{j-1},\mathbf i_j)$.
\item Let $\mathbf v_j=(\mathbf e_j)_3$.
\item Let $\mathbf w_j=(\mathbf e_j)_2$.
\end{enumerate}
The sample is then $\mathbf w=(\mathbf w_1,\ldots,\mathbf w_t)$.
We present two simple claims about $\mathbf W_{H,t}$.
\paragraph{Claim23.} For any $t\in\mathbb N$ and a $d$-regular hypergraph $H$ over $n$ vertices, $\mathbf W_{H,t}$ is $(nd^t)$-discretizable.
\begin{proof}
The sample $\mathbf w\sim\mathbf W_{H,t}$ is a deterministic function of $\mathbf v_0$ and $\mathbf i_1,\ldots,\mathbf i_t$, and those variables are set to a uniform choice out of $nd^t$ possibilities. The claim then follows from Definition5.
\end{proof}
\paragraph{Claim24.} For any $t\in\mathbb N$, $d$-regular hypergraph $H$, and $j\in[t]$, $\mathbf e_j$ is uniform over all $nd$ edges of $H$. As a result, $\mathbf W_{H,t}$ is homogeneous.
\begin{proof}
We will first prove that each $\mathbf e_j$ is uniform over $E_H$ by induction on$j$. $\mathbf e_1$ is uniform over the $nd$ edges in $E_H$ as it is sampled by independently choosing a starting vertex $\mathbf v_0\sim V$ uniformly and then uniformly choosing one of its $d$-neighbors. Furthermore, for any $j\in[t-1]$, if $\mathbf e_j$ is uniform, then $\mathbf v_j$ is also uniform. Hence, $\mathbf e_{j+1}$ is sampled by selecting a uniform vertex and then (independently) one of its $d$-neighbors, so $\mathbf e_{j+1}$ is also uniform over $E_H$.
Next, we show $\mathbf W_{H,t}$ is homogeneous. Fix any $a\in[n]$ and $j\in[t]$. As $\mathbf w_j=a$ if and only if $\mathbf e_j$ is one of the $d$-edges in $E_H$ whose second vertex is$a$ and $\mathbf e_j$ is uniform over the $nd$ edges in $E_H$, $\Pr[\mathbf w_j=a]=d/(nd)=1/n$.
\end{proof}
Finally, we show that whenever $H$ is a good expander, $\mathbf W_{H,t}$ is a good parity sampler.
\paragraph{Theorem25.} For any $n,d,t\in\mathbb N$, $\lambda,\varepsilon_0>0$, and $H$ a $\lambda$-spectral $d$-regular3-uniform hypergraph on$n$ vertices, $\mathbf W_{H,t}$ is an $(\varepsilon_0,\varepsilon\coloneqq(\varepsilon_0+\lambda)^t)$-parity sampler.
As an immediate corollary of Theorem25 and Proposition15:
\paragraph{Corollary26.} There exists an absolute constant $c_{\mathrm{spec}}>0$ for which the following holds. For any $n,d,t\in\mathbb N$, $\lambda,\varepsilon_0>0$, and $H$ a $\lambda$-mixing symmetric $d$-regular3-uniform hypergraph on$n$ vertices, $\mathbf W_{H,t}$ is an $(\varepsilon_0,\varepsilon\coloneqq(\varepsilon_0+c_{\mathrm{spec}}\cdot\lambda\log(1/\lambda))^t)$-parity sampler.
Throughout the remainder of this section, we will use the shorthand $\mathbf W\coloneqq\mathbf W_{H,t}$. For proving Theorem25, it will be convenient to consider $\sigma\in\{\pm1\}^n$ rather than $z\in\mathbb F_2^n$ (as in Definition9) using the mapping $\sigma_i=(-1)^{z_i}$. Equivalent to Definition9, $\mathbf W$ is an $(\varepsilon_0,\varepsilon)$ parity sampler if $|\bias_{\mathbf W}(\sigma)|\leq\varepsilon$ for all $\sigma\in\{\pm1\}^n$ satisfying $|\bias(\sigma)|=|\mathbb E_{i\sim[n]}[\sigma_i]|\leq\varepsilon_0$, where
\[\bias_{\mathbf W}(\sigma)\coloneqq\mathbb E_{\mathbf w\sim\mathbf W}\left[\prod_{j=1}^t\sigma_{\mathbf w_j}\right].\]
Similarly to[34], we express that bias algebraically.
\paragraph{Lemma27.} Let $A^{(\sigma)}\in\mathbb R^{n\times n}$ be defined as
\[A^{(\sigma)}_{i,k}\coloneqq\frac1d\sum_{(i',j',k')\in E_H}\sigma_{j'}\cdot1[i'=i\land k'=k].\]
Then, $\bias_{\mathbf W}(\sigma)=\frac1n\mathbf1^{\dagger}(A^{(\sigma)})^t\mathbf1$.
\begin{proof}
Let $\mathbf v_0,\ldots,\mathbf v_t$ and $\mathbf w_1,\ldots,\mathbf w_t$ be the random variables defined in the construction of $\mathbf W$. We claim that for each $j\in\{0,1,\ldots,t\}$ and$v\in[n]$, that
\[\mathbb E\left[1[\mathbf v_j=v]\prod_{k=1}^j\sigma_{\mathbf w_k}\right]=\frac1n\left(\mathbf1^{\dagger}(A^{(\sigma)})^j\right)_v\tag{6}\]
By induction on$j$. Clearly, for$j=0$, Equation(6) holds as both sides are equal to$1/n$ for any$v\in[n]$. For$j\geq1$,
\begin{align*}
&\mathbb E\left[1[\mathbf v_j=v]\prod_{k=1}^j\sigma_{\mathbf w_k}\right]\\
&=\sum_{v'\in[n]}\mathbb E\left[1[\mathbf v_j=v\land\mathbf v_{j-1}=v']\prod_{k=1}^j\sigma_{\mathbf w_k}\right]\\
&=\sum_{v'\in[n]}\mathbb E\left[1[\mathbf v_{j-1}=v']\prod_{k=1}^{j-1}\sigma_{\mathbf w_k}\right]\mathbb E[1[\mathbf v_j=v]\sigma_{\mathbf w_j}\mid\mathbf v_{j-1}=v']\\
&=\sum_{v'\in[n]}\frac1n\left(\mathbf1^{\dagger}(A^{(\sigma)})^{j-1}\right)_{v'}\mathbb E[1[\mathbf v_j=v]\sigma_{\mathbf w_j}\mid\mathbf v_{j-1}=v']\\
&=\frac1n\sum_{(a,b,c)\in E_H}\left(\mathbf1^{\dagger}(A^{(\sigma)})^{j-1}\right)_a\cdot\frac1d\cdot\sigma_b\cdot1[v=c]\\
&=\frac1n\left(\mathbf1^{\dagger}(A^{(\sigma)})^j\right)_v,
\end{align*}
where the third equality is the inductive hypothesis. Then,
\begin{align*}
\bias_{\mathbf W}(\sigma)&=\mathbb E_{\mathbf w\sim\mathbf W}\left[\prod_{j=1}^t\sigma_j\right]\\
&=\sum_{v\in[n]}\mathbb E\left[1[\mathbf v_j=v]\prod_{k=1}^j\sigma_{\mathbf w_k}\right]=\frac1n\mathbf1^{\dagger}(A^{(\sigma)})^t\mathbf1.
\end{align*}
\end{proof}
Next, we use the fact that$H$ is a$\lambda$-spectral expander to reason about$A^{(\sigma)}$.
\paragraph{Proposition28.} Let$J_n\in\mathbb R^{n\times n}$ be the matrix in which every element is$1/n$. For any$\sigma\in\{\pm1\}^n$,
\[\|A^{(\sigma)}-\bias(\sigma)J_n\|_{\mathrm{op}}\leq\lambda.\]
\begin{proof}
Fix any$x,z\in\mathbb R^n$. Then,
\begin{align*}
z^{\dagger}(A^{(\sigma)}-\bias(\sigma)J_n)x&=\frac1d\sum_{(i,j,k)\in E_H}x_k\sigma_jz_i-\frac{\bias(\sigma)}n\sum_{i\in[n]}x_i\sum_{j\in[n]}z_j\\
&=\frac1d\sum_{(i,j,k)\in E_H}x_i\sigma_jz_k-\frac1{n^2}\sum_{k\in[n]}\sigma_k\sum_{i\in[n]}x_i\sum_{j\in[n]}z_j\\
&\leq\lambda\|x\|_2\|z\|_2.\qquad\text{(Definition12)}
\end{align*}
\end{proof}
As an immediate consequence, we have:
\paragraph{Corollary29.} For any$\sigma\in\{\pm1\}^n$, $\|A^{(\sigma)}\|_{\mathrm{op}}\leq|\bias(\sigma)|+\lambda$.
\begin{proof}
As$\|J_n\|_{\mathrm{op}}=1$, the desired result follows from the reversed triangle inequality applied to the operator norm.
\end{proof}
Finally, we prove Theorem25.
\begin{proof}[Proof of Theorem25]
For any$\sigma\in\{\pm1\}^n$ satisfying$|\bias(\sigma)|\leq\varepsilon_0$, we have:
\begin{align*}
|\bias_{\mathbf W}(\sigma)|&=\frac{\mathbf1^{\dagger}}{\sqrt n}(A^{(\sigma)})^t\frac{\mathbf1}{\sqrt n}\qquad\text{(Lemma27)}\\
&\leq\|(A^{(\sigma)})^t\|_{\mathrm{op}}\leq\|A^{(\sigma)}\|_{\mathrm{op}}^t\\
&\leq(|\bias(\sigma)|+\lambda)^t\leq(\varepsilon_0+\lambda)^t.\qquad\text{(Corollary29)}
\end{align*}
\end{proof}
